\documentclass[11pt,a4paper]{article}
\usepackage[utf8]{inputenc}\usepackage[T1]{fontenc}
\usepackage[provide=*,english]{babel}
\usepackage{lmodern,microtype,amsmath,amssymb,booktabs,array,geometry,xcolor,fancyhdr}
\usepackage[hidelinks]{hyperref}
\geometry{margin=2.3cm,headheight=14pt}
\setlength{\parindent}{0pt}\setlength{\parskip}{0.4em}
\definecolor{gcblue}{RGB}{34,73,110}\definecolor{gcgray}{RGB}{90,96,102}
\pagestyle{fancy}\fancyhf{}
\fancyhead[L]{\small\color{gcgray}GC2D: BM4 implicit}
\fancyhead[R]{\small\color{gcgray}Model theory}\fancyfoot[C]{\thepage}
\newcommand{\R}{\mathbb{R}}
\begin{document}
\begin{center}
 {\LARGE\bfseries\color{gcblue}BM4 Implicit Method}\par
 \vspace{0.35em}{\large Physical state and one reduced Hairer projection per complete cycle}
\end{center}

\section{Twelve-stage BM4 base map}

Let $\chi_h$ be a first-order map and $\chi_h^*=\chi_{-h}^{-1}$ its
adjoint.  One complete step uses the palindromic composition
\[
\Psi_h=\chi_{b_{12}h}\circ\chi^*_{b_{11}h}\circ\cdots
\circ\chi_{b_2h}\circ\chi^*_{b_1h},
\]
with alternating adjoint and direct maps in chronological order and
\[
(b_1,\ldots,b_{12})=(a_1,a_2,a_3,a_4,a_5,a_6,a_6,a_5,a_4,a_3,a_2,a_1),
\]
\begin{align*}
a_1&=0.0792036964311957,&a_2&=0.1303114101821663,\\
a_3&=0.2228614958676077,&a_4&=-0.3667132690474257,\\
a_5&=0.3246481886897062,&a_6&=0.1096884778767498.
\end{align*}
The half-cycle coefficients sum to $1/2$.  Palindromy makes the composition
self-adjoint and the order conditions give designed global order four.  The
negative coefficient $a_4$ is an intentional backward subflow.

\section{Physical state and internal workspace}

The only accepted state is the physical guiding-centre state
$z\in\R^{2N}$.  Define
\[
E=\begin{pmatrix}I\\I\end{pmatrix},\qquad
P=\tfrac12\begin{pmatrix}I&I\end{pmatrix},\qquad
G=\begin{pmatrix}I&-I\end{pmatrix},\qquad
N=G^T.
\]
The doubled vector in $\R^{4N}$ is temporary algebraic workspace required by
the direct--adjoint splitting and Hairer projection.  It is neither an
accepted state nor a public state-extension option.  Time remains an external
integration parameter; no time-conjugate momentum is added to the state.

\section{One Hairer projection around the complete cycle}

For every complete BM4 step, the method displaces the diagonal embedding by a
single multiplier $\mu\in\R^{2N}$,
\begin{align}
\widehat Y_n&=Ez_n+N\mu,
&M(\mu)&=\Psi_{h,t_n}(\widehat Y_n),\\
r(\mu)&=GM(\mu)+2\mu,
&z_{n+1}&=P\bigl(M(\mu)+N\mu\bigr).
\end{align}
The accepted multiplier is the root $r(\mu)=0$.  Thus the projection is
applied once around the full twelve-stage composition, never between its
internal stages.

With $J_{\mathrm{BM4}}=D\Psi_{h,t_n}(\widehat Y_n)$, Newton uses
\[
D_\mu r=GJ_{\mathrm{BM4}}N+2I.
\]
The implementation forms $J_{\mathrm{BM4}}$ as the exact ordered product of
the twelve guiding-centre stage Jacobians, with a centred-difference fallback.
Good Broyden is retained as a nonlinear-solver choice; it changes the finite
iteration path but not the defining projected map.

\section{Single public method}

The model exposes only \texttt{BM4Implicit}.  Its numerical choices are
solver controls rather than distinct model classes:
\begin{center}
\small
\begin{tabular}{@{}ll@{}}
\toprule
Fixed architectural choice & Value\\
\midrule
Projection placement & once around each complete BM4 cycle\\
Projection formulation & reduced Hairer multiplier\\
Accepted state & physical $z\in\R^{2N}$\\
Temporary base-map workspace & doubled state in $\R^{4N}$\\
Public state extension & none\\
\bottomrule
\end{tabular}
\end{center}

The configurable controls are the coupling frequency, Newton tolerances,
maximum iterations, analytic or finite-difference Newton Jacobian, Newton or
Broyden nonlinear solver, progress reporting, and an optional step observer.

\section{Fixed-grid lifecycle and diagnostics}

The accepted physical state is advanced on the output-independent grid.  An
off-grid requested sample uses a shadow step from the preceding main node and
is discarded after sampling, so output density does not alter the accepted
trajectory.  Only main steps contribute observer records and diagnostics.

Each accepted step reports nonlinear iterations, residual evaluations, final
residual and tolerance, the projection-multiplier norm, coupling frequency,
and Jacobian strategy.  An optional \texttt{ImplicitBM4IntegrationStep}
contains the physical input/output, multiplier, solver metrics, and twelve
reconstructed base-stage snapshots.  Generalized energy may be reconstructed
as a diagnostic from those snapshots without introducing a public extended
state.

\section{Implementation correspondence}

The public implementation is \texttt{src/simulation/methods/bm4/implicit.py}.
The private \texttt{\_core.py} module owns only the coefficient schedule and
twelve-stage base-map traversal.  The detailed reduced-projection derivation
is retained in \texttt{implicit-reduced.tex}; runtime architecture is under
\texttt{../simulation/}.

\end{document}
